% runtime/runtime_services-DE.tex

\chapter{Laufzeitdienste}
\label{chap:runtime-services}

\index{Laufzeitdienste}
\index{Architektur!Laufzeit}

Laufzeitdienste sind die Ausführungsebene, welche genehmigtes
ingenieurtechnisches Wissen in operativen Kontexten bewertet. Sie verwenden
bestehende Objekte aus Grundlagen und Zustand und erzeugen abgeleitete
Laufzeitgrößen. Sie sind selbst keine Engineering Objects und definieren die
vorgelagerte Semantik nicht neu.

Dieses Kapitel beantwortet, wie Laufzeitbewertung organisiert wird, sodass
Projektion, Rahmenkonstruktion, dynamische Größen und Darstellungsgrenzen mit
der kanonischen AIM-Semantik konsistent bleiben. Es begründet die Architektur
der Laufzeitdienste und deren ingenieurtechnische Folgen für die nachgelagerte
Interpretation.

\section{Laufzeitinterpretation}

\begin{definition}[Laufzeitdienst]
\label{def:runtime-service}
\index{Laufzeitdienst!Definition}
Ein Laufzeitdienst ist ein operatorbasierter Bewertungsprozess, der während
der operativen Nutzung auf Alignment-Zuständen wirkt.
\end{definition}

\begin{principle}[Prinzip der Laufzeitbewertung]
\label{princ:runtime-evaluation}
Operative Größen sollen aus den zugrunde liegenden Zustandsbeschreibungen
bewertet und nicht redundant gespeichert werden.
\end{principle}

Laufzeitdienste bewerten Geometrie, Projektionen, Rahmen,
realisierungsabhängige Zustände, dynamische Größen und fahrzeugbezogene
Interpretation durch ausdrückliche Operatorkomposition.

Die folgenden Abschnitte folgen demselben Leserweg: von der Bewertung des
Laufzeitzustands über konkrete Dienste und die Komposition von Pipelines bis
zu Darstellungsgrenzen und architektonischen Folgen.

\section{Bewertung des Laufzeitzustands}

Die Laufzeitbewertung wird vom kanonischen Bewertungsoperator aus
\cref{def:evaluation-operator} bestimmt.

\begin{definition}[Laufzeitbewertung]
\label{def:runtime-evaluation}
\index{Laufzeitbewertung}
Die Laufzeitbewertung wird dargestellt durch
\[
\mathcal E
:
\mathcal X
\times
\Omega_{\mathrm{track}}
\rightarrow
\mathcal Y.
\]
\end{definition}

Dabei bezeichnet \(\mathcal X\) den Eisenbahnzustandsraum,
\(\Omega_{\mathrm{track}}\) den intrinsischen Gleisraum und \(\mathcal Y\)
die zur Laufzeit bewerteten Größen.

Typische Laufzeitausgaben umfassen World-Koordinaten, Tangenten- und
Rahmengrößen, aus Krümmung und Überhöhung abgeleitete Größen,
Projektionszustände und fahrzeugbezogen abgeleitete Größen.

\section{Persistente Information gegenüber abgeleitetem Laufzeitzustand}

In AIM speichern persistente Daten dünn besetzte, semantisch stabile
ingenieurtechnische Information. Laufzeitdienste rekonstruieren daraus
abgeleitete Größen.

\begin{principle}[Prinzip der Laufzeitrekonstruktion]
\label{princ:runtime-reconstruction}
\index{Prinzip der Laufzeitrekonstruktion}
Laufzeitgeometrie soll aus dünn besetzter Alignment-Semantik rekonstruiert und
nicht redundant gespeichert werden.
\end{principle}

Diese Trennung vermeidet Synchronisationsinstabilität, redundante Speicherung
und Inkonsistenzen der Realisierung.

\begin{figure}[ht]
\centering
% figures/runtime/support_runtime_to_representation_boundary-DE.tex
\begin{tikzpicture}[
  node distance=1.4cm,
  box/.style={draw, rounded corners, minimum width=4.7cm, minimum height=0.85cm, align=center},
  sec/.style={draw, rounded corners, dashed, minimum width=4.7cm, minimum height=0.85cm, align=center},
  arrow/.style={->, thick}
]
  \node[box] (rt) {Laufzeit};
  \node[sec, below=of rt] (rep) {Darstellung};
  \draw[arrow] (rt) -- (rep);
  \draw[thin] ($(rt.center)+(-3.0,-0.95)$) -- ($(rt.center)+(3.0,-0.95)$);
\end{tikzpicture}

\caption{Unterstützende Grenze: Die Laufzeitbewertung bleibt der Darstellung vorgelagert.}
\label{fig:support-runtime-representation-boundary}
\end{figure}

Die Abbildung macht die Verantwortungsgrenze ausdrücklich: Die Laufzeit
erzeugt abgeleiteten ingenieurtechnischen Zustand; die Darstellung verwendet
diesen Zustand zur Kodierung oder Visualisierung. Daraus folgt, dass
Laufzeitdienste anhand der Korrektheit ihrer Bewertung und nicht anhand
formatspezifischer Darstellungsentscheidungen validiert werden müssen.

\section{Laufzeitbewertung von pose3}

Die grundlegende pose3-Definition wird durch \cref{def:pose3} und
\cref{def:runtime-pose3} festgelegt. Laufzeitdienste verwenden diese
Begriffe und bewerten sie operativ.

\begin{definition}[Laufzeitbewertung von pose3]
\label{def:runtime-pose3-evaluation}
\index{Zustand!pose3!Laufzeitbewertung}
Die Laufzeitbewertung von pose3 kann schematisch dargestellt werden als
\[
\mathrm{pose3}_{\mathrm{run}}(s)
=
\mathcal E_{\mathrm{pose3}}(x,s).
\]
\end{definition}

Der pose3-Laufzeitzustand ist damit ein abgeleiteter Bewertungszustand und
keine persistente Engineering Identity.

\section{Projektionsdienste}

Projektionsdienste zur Laufzeit verwenden die kanonischen Operatoren gemäß
\cref{def:track-to-world-operator} und \cref{def:world-to-track-operator}.

Die folgenden beiden Abbildungen beantworten vor den Dienstdefinitionen eine
Richtungsfrage: Welcher Operator bildet von World nach Track ab, und welcher
liefert die komplementäre Abbildung von Track nach World?

\begin{figure}[ht]
\centering
\begin{tikzpicture}[
  node distance=1.25cm,
  box/.style={draw, rounded corners, minimum width=3.4cm, minimum height=0.8cm, align=center},
  op/.style={draw, rounded corners, minimum width=3.9cm, minimum height=0.85cm, align=center},
  arrow/.style={->, thick}
]
  \node[box] (world) {World Space};
  \node[op, below=of world] (opw) {World-to-Track-Operator};
  \node[box, below=of opw] (track) {Track Space};
  \draw[arrow] (world) -- (opw);
  \draw[arrow] (opw) -- (track);
\end{tikzpicture}

\caption{Unterstützende Operatorsicht: World--Track-Abbildung als gerichteter Laufzeitoperator.}
\label{fig:support-world-to-track-operator}
\end{figure}

\begin{figure}[ht]
\centering
\begin{tikzpicture}[
  node distance=1.25cm,
  box/.style={draw, rounded corners, minimum width=3.4cm, minimum height=0.8cm, align=center},
  op/.style={draw, rounded corners, minimum width=3.9cm, minimum height=0.85cm, align=center},
  arrow/.style={->, thick}
]
  \node[box] (track) {Track Space};
  \node[op, below=of track] (opt) {Track-to-World-Operator};
  \node[box, below=of opt] (world) {World Space};
  \draw[arrow] (track) -- (opt);
  \draw[arrow] (opt) -- (world);
\end{tikzpicture}

\caption{Unterstützende Operatorsicht: Track--World-Abbildung als komplementärer gerichteter Laufzeitoperator.}
\label{fig:support-track-to-world-operator}
\end{figure}

Gemeinsam begründet das Abbildungspaar vor den formalen Dienstdefinitionen
einen gerichteten Laufzeitvertrag: inverse und Vorwärtsprojektion sind
verwandte, aber nicht identische Operatoren. Daraus folgt die ausdrückliche
Trennung von numerischer Strategie, Mehrdeutigkeitsbehandlung und
Qualitätskriterien für jede Abbildungsrichtung.

\begin{definition}[Track--World-Laufzeitdienst]
\label{def:track-to-world-runtime}
\index{Projektion!Track--World-Laufzeitdienst}
Der Track--World-Laufzeitdienst bewertet
\[
\mathcal P_{\mathrm{tw}}
:
\Omega_{\mathrm{track}}
\rightarrow
\Omega_{\mathrm{world}}.
\]
\end{definition}

\begin{definition}[World--Track-Laufzeitdienst]
\label{def:world-to-track-runtime}
\index{Projektion!World--Track-Laufzeitdienst}
Der inverse Projektionsdienst zur Laufzeit bewertet
\[
\mathcal P_{\mathrm{wt}}
:
\Omega_{\mathrm{world}}
\rightarrow
\Omega_{\mathrm{track}}.
\]
\end{definition}

Der inverse Dienst ist nichtlinear, potenziell mehrdeutig, metrikabhängig und
kontextabhängig. Projektion bleibt von CRS-Transformation verschieden, wie in
\cref{def:coordinate-transform-operator} begründet.

Diese Unterscheidung ist die praktische Randbedingung für alle folgenden
Laufzeitdienste: Projektion interpretiert intrinsische Geometrie, während eine
CRS-Transformation nur die Darstellung im World-Raum ändert.

\section{Rahmendienste}

Die Rahmenkonstruktion zur Laufzeit verwendet den kanonischen Rahmenoperator
aus \cref{def:frame-operator}.

\begin{definition}[Rahmendienst]
\label{def:frame-service}
\index{Bezugsrahmen!Laufzeitdienst}
Der Rahmendienst bewertet
\[
\mathcal F(x_{\mathrm{metric}},s)
=
(\mathbf t(s),\mathbf n(s),\mathbf b(s)).
\]
\end{definition}

Rahmendienste unterstützen Überhöhungsinterpretation, pose3-Bewertung,
fahrzeugbezogene Interpretation, Projektionsinterpretation und
Laufzeitdynamik. Ihre Bewertung kann von der metrischen Realisierung abhängen.

\section{Abgeleitete dynamische Laufzeitgrößen}

Viele operative Größen werden zur Laufzeit abgeleitet und sind kein
persistentes Alignment-Wissen.

\begin{definition}[Effektive Querbeschleunigung]
\label{eq:runtime-effective-acceleration}
\index{Beschleunigung!effektive Querbeschleunigung}
\[
a_{\mathrm{eff}}
=
v^2\kappa-g\sin\phi.
\]
\end{definition}

Dies ist die Laufzeitform der kanonischen Beziehung aus
\cref{eq:effective-lateral-acceleration}.

\begin{definition}[Ruck zur Laufzeit]
\label{eq:runtime-jerk}
\index{Ruck!Laufzeit}
\[
j
=
v^3\kappa'
-
gv\cos\phi\,\phi'.
\]
\end{definition}

Dies ist die Laufzeitform der kanonischen Ruckbeziehung aus
\cref{eq:jerk-relation}.

\section{Hybride und LOD-abhängige Bewertung}

\begin{principle}[Prinzip hybrider Laufzeit]
\label{princ:hybrid-runtime}
Effiziente Laufzeitsysteme verbinden näherungsweise Grobbewertung,
zwischengespeicherte Darstellungen und exakte lokale Verfeinerung.
\end{principle}

\begin{definition}[LOD-abhängige Bewertung]
\label{def:lod-evaluation}
\index{Detaillierungsgrad (LOD)!Laufzeitbewertung}
Die Laufzeitbewertung hängt von Maßstab, Genauigkeitsanforderungen,
Sichtbarkeit, operativem Kontext, Interaktionszustand und Laufzeitkosten ab.
\end{definition}

Diese Verbindung ist zentral für skalierbare Projektion, Rahmenbewertung,
Visualisierung und fahrzeugbezogene Interpretation.

\section{Zwischenspeicherung und Konsistenz}

\begin{definition}[Laufzeitcache]
\label{def:runtime-cache}
\index{Laufzeitcache}
Ein Laufzeitcache speichert bereits bewertete Größen, um wiederholte
Berechnung zu vermeiden.
\end{definition}

Zwischenspeicherung kann Projektions-, Rahmen-, abgetastete Geometrie- und
dynamische Abfragen beschleunigen, führt aber Anforderungen an Konsistenz und
Invalidierung ein.

\section{Metrikbewusste Laufzeit}

Die Laufzeitbewertung hängt von der metrischen Realisierung ab
(\cref{chap:metric-realization}). Verschiedene Realisierungen beeinflussen die
Interpretation von Abständen, krümmungsbezogene Größen, Rahmenkonstruktion,
Projektionsverhalten und Stationierungsinterpretation.

\begin{principle}[Prinzip metrikbewusster Laufzeit]
\label{princ:metric-aware-runtime}
Laufzeitsysteme sollen Geometrie durch metrikbewusste Operatordienste und
nicht durch fest kodierte euklidische Annahmen bewerten.
\end{principle}

\section{Laufzeitpipelines und Zusammenwirken der Dienste}

\begin{definition}[Pipeline von Laufzeitoperatoren]
\label{def:runtime-pipeline}
\index{Pipeline!Laufzeitoperator}
Eine Laufzeitpipeline ist eine zusammengesetzte Operatorkette wie
\[
\mathcal E_{\mathrm{run}}
\sim
\mathcal P
\circ
\mathcal F
\circ
\mathcal G.
\]
\end{definition}

Die folgenden Abbildungen sind als zunehmend reichhaltige Antworten auf
dieselbe Frage zu lesen: Wie ändert sich die Operatorkette, wenn realisierte
Geometrie ausdrücklich einbezogen werden muss?

\begin{figure}[ht]
\centering
% figures/runtime/runtime_pipeline-DE.tex
\begin{tikzpicture}[
  node distance=0.85cm,
  box/.style={draw, rounded corners, align=center, minimum width=5.0cm, minimum height=0.7cm},
  arrow/.style={->, thick}
]
\node[box] (sparse) {Dünn besetztes Alignment};
\node[box, below=of sparse] (pose2) {pose2-Rekonstruktion};
\node[box, below=of pose2] (pose3) {pose3-Laufzeitzustand};
\node[box, below=of pose3] (dyn) {Laufzeitdynamik\\$\mathcal D$};
\node[box, below=of dyn] (vehicle) {Fahrzeuginterpretation\\$\mathcal V$};
\node[box, below=of vehicle] (opt) {Operative Bewertung\\$\mathcal O$};
\draw[arrow] (sparse) -- (pose2);
\draw[arrow] (pose2) -- (pose3);
\draw[arrow] (pose3) -- (dyn);
\draw[arrow] (dyn) -- (vehicle);
\draw[arrow] (vehicle) -- (opt);
\end{tikzpicture}

\caption{Begriffliche Pipeline der Laufzeitbewertung.}
\label{fig:runtime-pipeline}
\end{figure}

Diese Abbildung wird als grundlegende Operatorkomposition für die
Laufzeitbewertung von Zielzuständen interpretiert.

\begin{definition}[Realisierte Laufzeitpipeline]
\label{def:realized-runtime-pipeline}
\index{Pipeline!realisierungsbewusste Laufzeit}
Wenn realisierte Geometrie erforderlich ist, kann die Laufzeitbewertung
begrifflich abhängen von
\[
\mathcal E_{\mathrm{real}}
\sim
\mathcal P
\circ
\mathcal F
\circ
\mathcal G
\circ
\mathcal R.
\]
\end{definition}

\begin{figure}[ht]
\centering
% figures/runtime/realization_pipeline-DE.tex
\begin{tikzpicture}[
  node distance=0.9cm,
  box/.style={draw, rounded corners, align=center, minimum width=5.0cm, minimum height=0.75cm},
  arrow/.style={->, thick}
]
\node[box] (target) {Zielzustand\\$\mathrm{pose3}_{\mathrm{target}}$};
\node[box, below=of target] (realization) {Realisierungsoperator\\$\mathcal R_{\mathrm{pose3}}$};
\node[box, below=of realization] (real) {Realisierter Zustand\\$\mathrm{pose3}_{\mathrm{real}}$};
\node[box, below=of real] (measured) {Gemessener / rekonstruierter Zustand};
\node[box, below=of measured] (runtime) {Laufzeitbetrieb};
\draw[arrow] (target) -- (realization);
\draw[arrow] (realization) -- (real);
\draw[arrow] (real) -- (measured);
\draw[arrow] (measured) -- (runtime);
\end{tikzpicture}

\caption{Realisierungsbewusste Pipeline der Laufzeitbewertung.}
\label{fig:realization-pipeline}
\end{figure}

Die realisierungsbewusste Abbildung ergänzt demgegenüber ausdrücklich
\(\mathcal R\), um die Abhängigkeit vom realisierten Zustand darzustellen.
Daraus folgt, dass die Dienstwahl angeben muss, ob Ziel- oder realisierter
Zustand benötigt wird, weil diese Wahl die aktive Operatorkette ändert.

Verschiedene Laufzeitdienste aktivieren abhängig von Genauigkeit, Kontext und
Zweck unterschiedliche Operatorketten.

Dies führt zur abschließenden Grenzaussage: Laufzeitbewertung darf operativ
variieren, aber die begriffliche Verantwortung für ingenieurtechnische
Bedeutung nicht verändern.

\section{Grenzen zwischen Laufzeit, Realisierung und Darstellung}

Die Laufzeit muss Ziel-, realisierte, gemessene und temporäre
Interaktionszustände unterscheiden. Laufzeitbewertung und physische
Realisierung sind gekoppelt, aber begrifflich verschieden: Realisierung
transformiert Zielzustände; Laufzeit bewertet Zustände.

Die Darstellung bleibt nachgelagert. Laufzeitdienste stellen abgeleitete
Zustände und Größen bereit; Darstellungsdienste zeigen, kodieren oder tauschen
diese Ergebnisse für bestimmte Anwendungen aus.

\section{Begriffliche Laufzeitarchitektur}

\begin{figure}[ht]
\centering
\begin{tikzpicture}[
node distance=1.5cm and 2cm,
box/.style={draw,rounded corners,align=center,minimum width=3.6cm,minimum height=1.0cm},
arrow/.style={->, thick}
]
\node[box] (sparse) {Dünn besetztes Alignment\\Semantische Darstellung};
\node[box, below=of sparse] (metric) {Metrische Realisierung\\\(\mathcal G\)};
\node[box, below=of metric] (eval) {Laufzeitbewertung\\\(\mathcal E\)};
\node[box, below left=of eval] (proj) {Projektionsdienste\\\(\mathcal P\)};
\node[box, below=of eval] (frame) {Rahmendienste\\\(\mathcal F\)};
\node[box, below right=of eval] (dyn) {Dynamische Bewertung};
\node[box, below=2.2cm of eval] (vehicle) {Abgeleiteter Laufzeitzustand\\Anwendungsinterpretation};
\draw[arrow] (sparse) -- (metric);
\draw[arrow] (metric) -- (eval);
\draw[arrow] (eval) -- (proj);
\draw[arrow] (eval) -- (frame);
\draw[arrow] (eval) -- (dyn);
\draw[arrow] (proj) -- (vehicle);
\draw[arrow] (frame) -- (vehicle);
\draw[arrow] (dyn) -- (vehicle);
\end{tikzpicture}

\caption{Laufzeitdienste transformieren persistente ingenieurtechnische Information durch metrikbewusste Bewertung in abgeleiteten Laufzeitzustand für die nachgelagerte Interpretation.}
\label{fig:runtime-architecture}
\index{Abbildungen!Laufzeitarchitektur}
\end{figure}

Die Architekturabbildung fasst das Kapitelargument in ausführbarer Form
zusammen: Persistente Semantik speist metrikbewusste Bewertung; die Bewertung
verzweigt sich in Dienstfamilien, die im abgeleiteten Laufzeitzustand
zusammenlaufen. Daraus folgt ein stabiler Integrationsvertrag für die
nachgelagerten Kapitel zu Realität, Modellierung und Optimierung.

\section{Kanonische Interpretation}

\begin{proposition}
In AIM sind Laufzeitdienste operatorbasierte Bewertungssysteme, die auf
bestehenden Engineering- und Zustandsbegriffen wirken.
\end{proposition}

\begin{proposition}
Abgeleitete Laufzeitgrößen sind Bewertungsergebnisse und keine persistente
Engineering Identity.
\end{proposition}

\begin{proposition}
Projektion, Rahmenkonstruktion, metrische Realisierung und Laufzeitdynamik
bilden eine gekoppelte Dienstarchitektur.
\end{proposition}

\begin{proposition}
Die Darstellung bleibt der Laufzeitbewertung nachgelagert.
\end{proposition}

\section{Verbindung zu AIM}

\begin{summary}
AIM-Laufzeitdienste verwenden persistente ingenieurtechnische Information und
Begriffe der Zustandsebene und erzeugen durch operatorbasierte Bewertung
abgeleiteten ingenieurtechnischen Laufzeitzustand.

Diese Architektur hält die Semantik stabil, macht die Bewertung ausdrücklich,
wahrt Darstellungsgrenzen und ermöglicht anwendungsspezifische Interpretation,
ohne vorgelagertes Wissen neu zu definieren.

Der Laufzeitteil schließt daher mit einer präzisen Folge für die nächsten
Teile: Nachgelagerte Realität, Modellierung, Optimierung und Anwendungen dürfen
Laufzeitdienste spezialisieren, verwenden dafür aber diese Architektur, statt
sie neu zu definieren.
\end{summary}
